% §6.2 Sentence-BERT principal component channel, 12-city panel, with
% within-city-centered pooled PCA treatment definition.
% Numbers locked against:
%   results/replications/baur_pooled_pca/baur_pooled_pca_table.csv
%   results/replications/meta_regression_pooled.json (Baur source)

\subsection{Sentence-BERT principal component channel across twelve markets}
\label{sec:baur-12city}

The Baur, Rosenfelder, and Lutz channel \citep{baur2023analysis} treats
listing language as a $768$-dimensional vector from a frozen sentence-
BERT encoder and uses the leading principal component direction as a
scalar treatment in a partially-linear DML estimator. The original
study computes the PCA per market on the local listing distribution.
On a single-market analysis this is the canonical choice, but on a
twelve-market panel it produces a treatment whose direction is
identified only locally and only up to a sign, so per-market PC1
estimates can carry opposite signs across markets even when the
underlying language-price relationship is direction-consistent.

\paragraph{Pooled within-city-centered PCA.} To obtain a stable
treatment definition that is comparable across markets, we replace the
per-market PCAs with a single global principal-component direction
extracted from the pooled within-city covariance of the $n = 3{,}851$
stacked embeddings.  We center each city's embeddings at its own mean
before pooling, which removes between-city vocabulary differences from
the global covariance and produces a leading eigenvector identified as
the within-city common axis.  This is the standard pooling for cross-
population PCA discussed in \citet{flury1984}, equivalent in
construction to a Linear Discriminant Analysis pooled covariance.
We do not standardize individual embedding dimensions before pooling
because sentence-BERT outputs are L2-normalized by design and per-
dimension scaling would inflate the anisotropic noise structure
documented by \citet{ethayarajh2019contextual}.  We pin the sign of
the leading eigenvector by the positive-sum-of-loadings convention and
$z$-score the per-listing projection within each city, so the
estimand on each market remains "per per-city-$\sigma$ move along the
common axis".

The leading global axis explains $6.1\%$ of the total variance in
the pooled embedding matrix, consistent with the anisotropic
clustering of sentence-BERT representations documented by
\citet{ethayarajh2019contextual}. Per-listing scores along this axis
correlate positively with utilitarian and starter-home vocabulary
("investor special", "as-is", "cosmetic", "TLC") and negatively with
luxury and amenity-rich vocabulary ("renovated", "chef's kitchen",
"designer finishes"); the sign of the per-$\sigma$ DML estimate
should consequently be negative if listings whose language sits
toward the utilitarian end of the axis tend to transact at lower
prices, which is the expected pattern in our hedonic pricing setup.

\paragraph{Per-market estimates.}
Table~\ref{tab:baur-12city} reports the per-market per-$\sigma$ effect
for each market under the pooled-PCA treatment.  Eleven of the twelve
markets return a negative point estimate, and the twelfth (Atlanta)
returns a near-zero positive estimate of $+0.006$. The sign-flip
pattern that the per-market PCA produced (four significant markets
with three negative and one positive sign) is absent: the corrected
treatment definition delivers a single direction of effect across the
panel.  Among the twelve markets, only Portland returns a confidence
interval that excludes zero at the $95\%$ level
($\hat\theta = -0.109$, CI $[-0.171, -0.052]$); the remaining eleven
have absolute point estimates between $0.001$ and $0.259$ with
confidence intervals that span zero.  Seattle has the largest absolute
point estimate of $-0.259$ but the widest standard error
($\hat{\mathrm{SE}} = 0.122$) and consequently the widest confidence
interval $[-0.384, +0.193]$.

\paragraph{Pooled inference.} Section~\ref{sec:meta-regression-12}
takes up the pooled inference question directly via random-effects
meta-analysis of the twelve per-market estimates. The Paule-Mandel +
HKSJ pooled per-$\sigma$ estimate is $-0.029$ with $95\%$ confidence
interval $[-0.051, -0.008]$ and $p = 0.013$, indicating a small but
statistically detectable negative effect across the twelve-metropolitan
panel.  We attempted a single pooled DML estimator with city fixed
effects and the Chiang-Kato-Ma-Sasaki \citep{chiang2022multiway}
cluster-robust variance, but the small number of clusters
($G = 12$) produced fold-assignment instability in the cross-fitted
ridge nuisances that we judged not robust enough for a publication
headline: across three independent invocations with the same seed and
data we observed pooled point estimates ranging from $-0.19$ to
$-0.89$. The random-effects meta-analysis treats each market's
$\hat\theta_i$ as a sufficient statistic and combines them via inverse-
variance weighting with a Paule-Mandel $\tau^2$ estimate, which is
the canonical small-$G$ alternative to cluster-respecting cross-
fitting \citep{veroniki2016methods}.

\paragraph{Dollarization.} Translating the pooled per-$\sigma$ effect
to a per-swap dollar amount at the median sale price across the
twelve metros gives an approximate $-2.9\%$ price decrement per
$\sigma$ move along the global axis toward the utilitarian end, or
roughly $-\$18{,}000$ at the panel-weighted-mean median sale price of
$\$636{,}000$. We caution that this is a panel-weighted average
across markets with substantially different median prices; the per-
market translations span $-\$31{,}000$ in San Francisco to $-\$10{,}000$
in Philadelphia.
